
% === Глава диссертации ===
\chapter{Оптимизация SQL-запросов на основе обратной связи с помощью больших языковых моделей (FaLQON)}
\label{ch:falqon}

\section*{Аннотация}
Оптимизация выполнения запросов является одной из важнейших задач в современных реляционных СУБД. Встроенные планировщики, основанные на стоимостных моделях, нередко генерируют субоптимальные планы. Использование подсказок (hints) позволяет направлять планировщик, но их ручная настройка требует высокой экспертизы. В главе предлагается \emph{Feedback-Driven LLM Query Optimizer Network (FaLQON)}~— итеративный агент на базе большой языковой модели (БЯМ), который генерирует наборы подсказок, получает обратную связь от СУБД (план, стоимость, фактическое время исполнения) и в диалоге улучшает стратегию. Экспериментальная оценка на бенчмарке JOB показывает снижение задержки исполнения до 28{.}5\% по сравнению с дефолтным планировщиком openGauss.

\bigskip
\noindent\textbf{Ключевые слова:} оптимизация запросов, большие языковые модели, подсказки планировщику, итеративный метод, обратная связь, агентная система.

\section{Введение}
СУБД играют ключевую роль в современных цифровых системах, а их производительность напрямую влияет на скорость работы приложений и бизнес-процессов. Фокус данной главы~— оптимизация работы планировщика запросов: компонента, который преобразует SQL-запрос в последовательность операций с минимальными издержками. Ошибки оценок (кардинальности, стоимости) и неполные/устаревшие статистики приводят к субоптимальным планам.

Одним из перспективных направлений является применение больших языковых моделей (БЯМ) в качестве \emph{внешнего советчика} планировщику: БЯМ умеют анализировать текст запроса, план и исторические метрики, предлагать подсказки (включение/отключение операторов, выбор стратегий соединения, и~т.\,п.). В этой главе мы формализуем задачу и описываем метод FaLQON~— интерактивный агент, который под действием обратной связи от СУБД уточняет подсказки и ускоряет запросы.

\section{Связанные работы}
\subsection{Глубокие нейронные сети для оптимизации запросов}
\textbf{BAO}~\cite{marcus2022bao}~— надстройка над классическим оптимизатором: базовый оптимизатор генерирует несколько планов, Tree-CNN~\cite{roy2019treecnn} предсказывает задержки, а выбор плана выполняется методом Томпсона~\cite{thompson1933}. \textbf{COOOL}~\cite{xu2023coool} переходит от точного предсказания стоимости к задаче ранжирования планов. \textbf{Robust-MSCN}~\cite{negi2023robust} повышает устойчивость оценки кардинальности под смену рабочих нагрузок за счёт маскирования признаков и битовых карт соединений.

\subsection{Большие языковые модели для оптимизации запросов}
Современные подходы включают \textbf{LLM-R2}~\cite{li2024llmr2} (переписывание запросов по правилам), \textbf{LLM-QO}~\cite{tan2025can} (генерация текстового представления плана БЯМ) и \textbf{LLM-Steer}~\cite{akioyamen2024unreasonable} (использование эмбеддингов SQL и классификатора подсказок). Новизна FaLQON состоит в организации \emph{диалога} между БЯМ и СУБД с итеративной оценкой качества подсказок по фактическим метрикам.

\section{Постановка задачи}
Пусть $\tau$~— исходная задача (например, SQL-запрос). Обозначим через $S(\tau; H)$ план/решение внешней системы при применении набора подсказок $H$, а через $S(\tau;\varnothing)$~— решение без подсказок. Введём функцию стоимости (например, время исполнения) $C(\cdot)$. Цель~— найти
\begin{equation}
  H^\star \;=\; \arg\min_{H \in \mathcal{H}(\tau)} \; C\!\bigl(S(\tau; H)\bigr),
  \label{eq:goal}
\end{equation}
где $\mathcal{H}(\tau)$~— множество допустимых комбинаций подсказок для задачи $\tau$. Поскольку $|\mathcal{H}(\tau)|$ экспоненциально велико, прямой перебор нереализуем.

При использовании БЯМ имеем вектор параметров $\theta$ и текстовое представление задачи $\tau$. Рекомендации формируются как
\begin{equation}
  H \;=\; \mathcal{G}_\theta\!\bigl(\tau, \iota\bigr),
  \label{eq:llm_one_shot}
\end{equation}
где $\iota$~— инструкция/шаблон промпта. На практике мы используем итеративную схему с обратной связью.

\section{Метод FaLQON}
\subsection{Интерактивная генерация подсказок}
Пусть $T(H)$~— фактическое время выполнения запроса при подсказках $H$, $T(\varnothing)$~— при настройках по умолчанию. Определим \emph{процентную награду}
\begin{equation}
  r(H) \;=\; 100 \cdot \frac{T(\varnothing) - T(H)}{T(\varnothing)}.
  \label{eq:reward}
\end{equation}
Разбивая значения $r(H)$ на интервалы, получаем категориальную награду $\mathcal{R}(H) \in \{0,1,\dots\}$. Тогда итеративная генерация имеет вид
\begin{equation}
  H^{(k+1)} \;=\; \mathcal{G}_\theta\!\bigl(\tau,\; \mathcal{D}_k,\; \mathcal{R}\!\bigl(H^{(k)}\bigr) \bigr), 
  \qquad k=0,1,2,\dots,
  \label{eq:iter_text}
\end{equation}
где $\mathcal{D}_k$~— накопленная история (запрос, метрики, контекст). В варианте с использованием текущего плана $P^{(k)} = S(\tau; H^{(k)})$:
\begin{equation}
  H^{(k+1)} \;=\; \mathcal{G}_\theta\!\bigl(\tau,\; P^{(k)},\; \mathcal{D}_k,\; \mathcal{R}\!\bigl(H^{(k)}\bigr) \bigr).
  \label{eq:iter_plan}
\end{equation}

\begin{figure}[t]
  \centering
  \includegraphics[width=.9\linewidth]{figures/fig1-hint-tree}
  \caption{Дерево комбинаций подсказок и их влияние на производительность.}
  \label{fig:hint-tree}
\end{figure}

\begin{figure}[t]
  \centering
  \includegraphics[width=.9\linewidth]{figures/fig2-falqon-arch}
  \caption{Общая схема архитектуры системы FaLQON.}
  \label{fig:falqon-arch}
\end{figure}

\subsection{Инструментарий и подготовка диалогов для обучения}
Пространство подсказок формируется полным перебором грубых подсказок для запросов из JOB, с организацией результатов в дерево (рис.~\ref{fig:hint-tree}). Набор поддерживаемых подсказок ограничен для контроля размера комбинаторики и задержек генерации. Пример списка:
\begin{itemize}
  \item \texttt{set(enable\_seqscan, off);} 
  \item \texttt{set(enable\_bitmapscan, off);} 
  \item \texttt{set(enable\_indexscan, off);} 
  \item \texttt{set(enable\_indexonlyscan, off);} 
  \item \texttt{set(enable\_hashjoin, off);} 
  \item \texttt{set(enable\_nestloop, off);} 
  \item \texttt{set(enable\_mergejoin, off).}
\end{itemize}

\begin{figure}[t]
  \centering
  \includegraphics[width=.9\linewidth]{figures/fig3-pipeline}
  \caption{Процесс получения и оценки обучаемой модели для корпуса SQL-запросов.}
  \label{fig:pipeline}
\end{figure}

\section{Экспериментальная оценка}
\subsection{Методология и стенд}
Эксперименты выполнены на одноузловом сервере (128 ядер HiSilicon Kunpeng~920, 1{\,}000~ГБ ОЗУ) с СУБД \emph{openGauss}~\cite{li2021opengauss}. Ключевые параметры:
\begin{quote}
\texttt{shared\_buffers = 256\,GB; \quad bulk\_write\_ring\_size = 10\,GB; \quad work\_mem = 80\,GB;\\
cstore\_buffers = 4\,GB; \quad query\_dop = 1.}
\end{quote}
В качестве БЯМ использована \emph{DeepSeek Coder 1.3B}~\cite{guo2024deepseek} (окно 16k токенов). Данные: JOB~\cite{leis2015howgood} (33 шаблона, соединения 4–17 таблиц). Часть шаблонов исключена из обучения для проверки обобщения.

\subsection{Результаты}
Обучающее множество диалогов по JOB составило 4{,}626 экземпляров (без четырёх случайных шаблонов). На рис.~\ref{fig:accel-a}--\ref{fig:accel-b} показано отношение достигнутого ускорения к максимальному для двух режимов (обучение на тексте запроса и на плане~запроса).
\begin{figure}[t]
  \centering
  \includegraphics[width=.9\linewidth]{figures/fig4-accel-text}
  \caption{Отношение достигнутого ускорения к максимальному: обучение на текстах запросов (а).}
  \label{fig:accel-a}
\end{figure}

\begin{figure}[t]
  \centering
  \includegraphics[width=.9\linewidth]{figures/fig5-accel-plan}
  \caption{Отношение достигнутого ускорения к максимальному: обучение на планах запросов (б).}
  \label{fig:accel-b}
\end{figure}

Подробные статистики времени выполнения приведены в табл.~\ref{tab:job-runtime}. Согласно данным, максимальное время выполнения удалось сократить в 1{.}94 и 2{.}24 раза за счёт подсказок БЯМ; результаты составляют 79{.}45\% и 88{.}62\% от условного оптимума.

\begin{table}[t]
  \centering
  \caption{Статистики времени выполнения запросов для JOB.}
  \label{tab:job-runtime}
  \begin{tabular}{llcccc}
    \toprule
    \textbf{Метрика} & \textbf{Набор} & \textbf{Оптимальный} & \textbf{БЯМ} & \textbf{Жадный} & \textbf{Без подсказок}\\
    \midrule
    Медиана, мс & Полный   & 8.94 & 9.41 & 9.28 & 1.36 \\
                & Тестовый & 1.30 & 1.34 & 1.35 & 1.51 \\
    \midrule
    Среднее, мс & Полный   & 3.19 & 3.55 & 3.31 & 4.97 \\
                & Тестовый & 3.40 & 3.67 & 3.61 & 5.74 \\
    \midrule
    Максимум, мс& Полный   & 3.45 & 3.67 & 3.45 & 7.17 \\
                & Тестовый & 1.35 & 1.36 & 1.36 & 3.05 \\
    \midrule
    Суммарное время, мс & Полный & 3.60 & 4.01 & 3.77 & 5.62 \\
                        & Тестовый & 5.10 & 5.50 & 5.42 & 5.62 \\
    \midrule
    Ускорение, \% & Полный   & 35.87 & 28.50 & 33.23 & 0 \\
                  & Тестовый & 40.78 & 36.14 & 37.14 & 0 \\
    \bottomrule
  \end{tabular}
\end{table}

Статистика по \emph{числу шагов} до достижения целевого ускорения приведена в табл.~\ref{tab:job-steps}. Решение на основе БЯМ близко к оптимальному и даёт $\sim\!5\times$ сокращение шагов по сравнению с жадным поиском.

\begin{table}[t]
  \centering
  \caption{Шаги до достижения целевого ускорения на JOB.}
  \label{tab:job-steps}
  \begin{tabular}{llccc}
    \toprule
    \textbf{Метрика} & \textbf{Набор} & \textbf{Оптимальный} & \textbf{БЯМ} & \textbf{Жадный} \\
    \midrule
    Медиана        & Полный   & 0 & 1 & 0 \\
                   & Тестовый & 0 & 1 & 0 \\
    \midrule
    Среднее        & Полный   & 1.96 & 1.71 & 8.56 \\
                   & Тестовый & 2.20 & 1.53 & 8.40 \\
    \midrule
    Максимум       & Полный   & 6 & 5 & 28 \\
                   & Тестовый & 5 & 5 & 21 \\
    \midrule
    Суммарное кол-во & Полный & 221 & 193 & 1001 \\
                     & Тестовый & 33 & 23 & 126 \\
    \bottomrule
  \end{tabular}
\end{table}

\section{Заключение}
Предложен и апробирован \emph{FaLQON}~— итеративный агент на базе БЯМ для автоматической оптимизации SQL-запросов. Ключевая особенность~— диалог с СУБД по фактическим метрикам исполнения. На JOB достигнуто снижение задержки на 28{.}5\% относительно openGauss по умолчанию при существенном сокращении числа шагов до целевого ускорения.

\section*{Благодарности}
Авторы благодарят сообщества openGauss и DeepSeek за открытые инструменты и данные.

\section*{Список литературы}
\begin{thebibliography}{99}
\bibitem{marcus2022bao} R.~Marcus, P.~Negi, H.~Mao, N.~Tatbul, M.~Alizadeh, and T.~Kraska. Bao. \emph{ACM SIGMOD Record}, 51(1):6--13, 2022.
\bibitem{roy2019treecnn} D.~Roy, P.~Panda, and K.~Roy. Tree-CNN: A hierarchical Deep Convolutional Neural Network for incremental learning. \emph{Neural Networks}, 121:148--160, 2019.
\bibitem{thompson1933} W.~R.~Thompson. On the likelihood that one unknown probability exceeds another in view of the evidence of two samples. \emph{Biometrika}, 25(3/4):285, 1933.
\bibitem{xu2023coool} X.~Xu, Z.~Zhao, T.~Zhang, R.~Kang, L.~Sun, and J.~Chen. COOOL: A Learning-To-Rank Approach for SQL hint Recommendations. \emph{arXiv}, 2023.
\bibitem{negi2023robust} P.~Negi et al. Robust Query Driven Cardinality Estimation under Changing Workloads. \emph{PVLDB}, 16(6):1520--1533, 2023.
\bibitem{li2024llmr2} Z.~Li, H.~Yuan, H.~Wang, G.~Cong, and L.~Bing. LLM-R2: a large language model enhanced rule-based rewrite system for boosting query efficiency. \emph{arXiv}, 2024.
\bibitem{tan2025can} J.~Tan et al. Can large language models be query optimizer for relational databases? \emph{arXiv}, 2025.
\bibitem{akioyamen2024unreasonable} P.~Akioyamen, Z.~Yi, and R.~Marcus. The unreasonable effectiveness of LLMs for query optimization. \emph{arXiv}, 2024.
\bibitem{li2021opengauss} G.~Li et al. OpenGauss. \emph{PVLDB}, 14(12):3028--3042, 2021.
\bibitem{guo2024deepseek} D.~Guo et al. DeepSeek-Coder. \emph{arXiv}, 2024.
\bibitem{leis2015howgood} V.~Leis et al. How good are query optimizers, really? \emph{PVLDB}, 9(3):204--215, 2015.
\end{thebibliography}

% === Примечание по рисункам ===
% Поместите файлы изображений в папку ./figures/ с именами:
%   fig1-hint-tree.(pdf|png|jpg)
%   fig2-falqon-arch.(pdf|png|jpg)
%   fig3-pipeline.(pdf|png|jpg)
%   fig4-accel-text.(pdf|png|jpg)
%   fig5-accel-plan.(pdf|png|jpg)
% Таблицы и формулы воспроизведены согласно исходной работе.
